%----------
%	CONCLUSION OF THESIS
%---------------------------------------------------------------------------------

\addchaptertocentry{Дүгнэлт}
\pagecolor{white}

\section*{\centering \huge{Дүгнэлт}}

Энэхүү дипломын ажлын хүрээнд хувь хүний санхүүгийн мэдээллийг нэгтгэн удирдах, автоматаар шинжлэх, хиймэл оюунд тулгуурласан зөвлөмж өгөх \textbf{FinTrack} системийг боловсруулж, судалгааны зорилго, дэвшүүлсэн зорилтуудыг хэрэгжүүллээ. Ажлын явцад дараах гол үр дүнгүүдэд хүрсэн болно.

\subsection*{Хүрсэн үр дүнгүүд}

\textbf{1. Шаардлагын инженерчлэлийн талаас.} Хэрэглэгч, админ гэсэн хоёр төрлийн оролцогчид 28 функциональ, 10 функциональ бус шаардлагыг боловсрууллаа. 20 хэрэглээний тохиолдлыг диаграмм болон дэлгэрэнгүй тодорхойлолтоор баримтжуулсан.

\textbf{2. Системийн зохиомжийн талаас.} Сонгодог гурван давхаргат архитектурыг өргөтгөж, AI давхаргатай дөрвөн давхаргат архитектурыг боловсруулсан. 11 хүснэгттэй ER загвар, 12 боловсруулагч файл, 45 REST API цэг, 17 React Native дэлгэц, AI шинжилгээний долоон алхамт хоолой зэргийг бүрэн зохиомжилсон.

\textbf{3. Хэрэгжүүлэлтийн талаас.} Сервер тал нь Go 1.25 + Gin + GORM, хэрэглэгч тал нь React Native + Expo SDK 55 + NativeWind v4, өгөгдлийн сан нь PostgreSQL 16, AI нь Gemini 2.5 Flash + OpenRouter нөөц гэсэн орчин үеийн технологийн багцаар бүтээсэн. Контейнержуулалт (Docker Compose) ба олон AI нийлүүлэгчтэй автомат нөөц шилжилтийн механизмыг хэрэгжүүлсэн.

\textbf{4. Туршилтын талаас.} 20 функциональ шаардлагаас 85\% нь бүрэн биелсэн. 30 бодит банкны хуулга дээр AI шинжилгээ F1 = 0.967 нарийвчлал үзүүлсэн. 8 хэрэглэгчтэй UAT туршилтад дундаж амжилтын хувь 93.8\%, SUS оноо 81.3 буюу "Маш сайн" түвшинд гарсан. OWASP Top-10-ийн 9 ангилалд бүгдийг өнгөрсөн.

\subsection*{Шинэлэг тал}

Энэхүү ажлын шинэлэг талыг доорх дөрвөн чиглэлээр тодорхойлж байна:

\textbf{1. Монгол хэлний дэмжлэгтэй LLM AI чат.} Mint, YNAB, Monarch Money зэрэг олон улсын ижил төстэй системүүд монгол хэлийг дэмждэггүй. FinTrack нь Gemini 2.5-ийн монгол хэлний чадварыг ашиглан хэрэглэгчтэй чөлөөтэй харилцах AI чатыг хэрэгжүүлсэн.

\textbf{2. Монголын банкны хуулгын автомат AI задлал.} ХХБ, Голомт, Хаан банкны PDF/Excel хуулгын тусгай форматыг таних Python задлагч ба LLM-д суурилсан ангилал, дүн шинжилгээний хосолсон хоолойг бүтээсэн. Энэ нь дэлхийн зах зээл дэх ямар ч ижил төстэй системийн дэмждэггүй өвөрмөц шинж чанар юм.

\textbf{3. Олон нийлүүлэгчтэй AI архитектур (газарзүйн хязгаарлалтыг тойрох).} Google Gemini API нь Монгол Улсаас шууд хандах боломжгүй тул автоматаар OpenRouter руу шилжих нөөц шилжилтийн стратегийг хэрэгжүүлж, үнэгүй загваруудын жагсаалтыг 1 цаг тутамд автомат шинэчилдэг механизмыг бий болгосон. Энэ нь хөгжиж буй орнуудад LLM-д суурилсан үйлчилгээ бүтээх практик асуудлыг шийдсэн арга зүйн жишээ юм.

\textbf{4. Контекстэд тулгуурласан санхүүгийн зааварчилгаа.} Хэрэглэгчийн сүүлийн 30 хоногийн гүйлгээ, нийт үлдэгдэл, ангиллын задаргааг шууд системийн зааварчилгаанд JSON хэлбэрээр оруулдаг бөгөөд энэ нь дахин сургалт шаарддаггүй, контекст дотор суралцах хандлагад тулгуурласан хувийн санхүүгийн зөвлөгөөний зарчмыг бий болгож байна.

\subsection*{Практик ач холбогдол}

Судалгааны практик ач холбогдол нь дараах талуудаар илэрнэ:
\begin{itemize}
    \item \textbf{Хэрэглэгчийн талаас} --- хувийн санхүүгийн сахилга батыг сайжруулах, төлөвлөлтийн чанарыг нэмэгдүүлэх, санхүүгийн мэдээллийг ойлгомжтой хэлбэрээр хүлээн авах боломж бүрдсэн. UAT-ын үр дүнгээр дундаж хэрэглэгч банкны хуулгаа гар утсаар $\sim$1 минутын дотор ангилуулж зөвлөгөө авах боломжтой болсон.
    \item \textbf{Үндэсний хэмжээнд} --- Монголын банкны хэрэглэгчдэд тохирсон, орчуулга шаардахгүй, локал контекст ойлгодог санхүүгийн AI шийдлийн анхны эх загвар бий болов.
    \item \textbf{Технологийн талаас} --- Монгол хөгжүүлэгчдэд гадны geo-block-той AI үйлчилгээг найдвартай ашиглах архитектурын жишээ үлдсэн.
    \item \textbf{Боловсролын талаас} --- Орчин үеийн full-stack хөгжүүлэлтийн (Go, React Native, PostgreSQL, LLM) практик хэрэгжүүлэлтийн нэгдсэн жишээ материал болсон.
\end{itemize}

\subsection*{Цаашдын хөгжүүлэлтийн чиглэл}

Системийг цаашид доорх чиглэлээр өргөтгөн хөгжүүлэх боломжтой:
\begin{enumerate}
    \item \textbf{Open Banking интеграци.} Монгол банкны зохицуулалт нь Open Banking-ыг бүрэн нэвтрүүлбэл хуулга гараар оруулахын оронд бодит цагийн нэгтгэл хийх боломжтой болно.
    \item \textbf{Хадгаламж, хөрөнгө оруулалтын модуль} (FR-16). Зорилт тавьж хадгаламж хянах, хувьцаа, бонд, крипто хөрөнгийн дансыг нэгтгэх.
    \item \textbf{Гэр бүлийн хамтын данс.} Олон хэрэглэгчтэй хуваалцах данс, зөвшөөрлийн систем.
    \item \textbf{Төхөөрөмж дээр ажиллах LLM.} Хувийн нууцыг хамгаалах үүднээс Gemma 2B/Llama 3.2 1B-ийг утсан дээр шууд ажиллуулах боломжийг судлах.
    \item \textbf{Таамаглалын загвар.} Time-series загварыг (Prophet, NeuralProphet) ашиглан ирээдүйн мөнгөн урсгал ба сарын зарлагын таамаглал гаргах.
    \item \textbf{Олон хэрэглэгчийн өргөтгөл.} 1{,}000+ зэрэг хэрэглэгчтэй ажиллах хэвтээ өргөтгөл, Redis кэш, Kubernetes байршуулалтын туршилтыг гүйцэтгэх.
    \item \textbf{Хуулга уншилтын нарийвчлал.} PDF-ийн зарим багана буруу таних асуудлыг шийдвэрлэхэд зориулсан банк бүрийн тусгай задлагч бичих.
\end{enumerate}

\subsection*{Эцсийн дүгнэлт}

Энэхүү ажил нь хувийн санхүүгийн AI шийдлийн салбарт Монгол хэрэглэгчдэд зориулсан анхны цогц эх загварыг бүтээж, дэлхийн ижил төстэй шийдлээс ялгарах гурван өвөрмөц шинж чанарыг (монгол хэл, монголын банкны хуулга задлал, газарзүйн хязгаарлалтыг тойрох олон нийлүүлэгчтэй архитектур) баталгаатай хэрэгжүүлсэн. Туршилтын тоон үзүүлэлтүүд нь системийн чанарыг бодит хэрэглэгчийн орчинд нотолсон бөгөөд цаашид үндэсний хэмжээний хувийн санхүүгийн ухаалаг туслахын суурь шийдэл болох боломжтойг харуулж байна.
